\documentclass[11pt]{article}

\usepackage{amsmath,amssymb}
\usepackage{physics}
\usepackage[utf8]{inputenc}
\usepackage[letterpaper,margin=1.2in]{geometry}
\usepackage[colorlinks=true,allcolors=blue]{hyperref}
\usepackage[makeroom]{cancel}
\usepackage{xcolor}
\usepackage{mdframed}

%% Color boxes for key results
\newmdenv[backgroundcolor=blue!8,linecolor=blue!40,linewidth=1pt,
          innertopmargin=6pt,innerbottommargin=6pt]{keyresult}

%% Notation (matching paper 2 notation.tex)
\newcommand{\iunit}{\imath}          % imaginary unit (dotless i)
\newcommand{\Avec}{\mathbf{A}}
\newcommand{\Evec}{\mathbf{E}}
\newcommand{\Bvec}{\mathbf{B}}
\newcommand{\phat}{\hat{\mathbf{p}}}
\newcommand{\rhat}{\hat{\mathbf{r}}}
\newcommand{\Vnl}{V^{\mathrm{nl}}}
\newcommand{\HLG}{\hat{H}^{\mathrm{LG}}}
\newcommand{\HVG}{\hat{H}^{\mathrm{VG}}}

%% Margin notes for teaching points
\newcommand{\teachnote}[1]{\marginpar{\footnotesize\color{blue!70}#1}}

\title{%
  Electromagnetic Interaction in RT-TDDFT:\\
  From the Plane Wave to the Long-Wavelength Approximation\\[4pt]
  {\large Pedagogical Derivation — Length and Velocity Gauge}%
}
\author{Notes for RMG Paper 2 (Jacek Jakowski)}
\date{April 2026}

\begin{document}
\maketitle
\tableofcontents
\newpage

%%=================================================================
\section{Purpose and Conventions}
%%=================================================================

These notes derive the electromagnetic interaction Hamiltonian
used in real-time TDDFT (RT-TDDFT) starting from the most general
description --- a monochromatic electromagnetic plane wave ---
and arriving step by step at the spatially uniform vector potential
$\Avec(t)$ used in the velocity-gauge implementation.
All intermediate steps are kept explicitly.
The goal is to make clear:
\begin{itemize}
  \item[(1)] what physical approximation the long-wavelength limit
    represents,
  \item[(2)] how the length and velocity gauges are related,
  \item[(3)] what new physics becomes accessible if the
    long-wavelength approximation is relaxed, and
  \item[(4)] how the Coulomb-matrix perturbation of
    ref.~[Jakowski et al.\ JCTC 2019] relates to the plane-wave
    framework.
\end{itemize}

\noindent\textbf{Units.}
All equations use \textbf{atomic units} unless stated otherwise:
$\hbar=m_e=e=4\pi\varepsilon_0=1$.
The speed of light is $c=1/\alpha\approx137$ (kept explicit for
clarity).
The electron has charge $-1$ (i.e.\ $q_e = -e = -1$).
In atomic units the minimal-coupling substitution
$\hat{\mathbf{p}}\to\hat{\mathbf{p}}-q_e\Avec/c
= \hat{\mathbf{p}}+\Avec$ (since $q_e=-1$, $c\to1$ absorbed into
$\Avec$ in the convention used here).

\noindent\textbf{Imaginary unit.}
We use $\iunit$ (dotless $i$) throughout, consistent with paper~1
and paper~2.

%%=================================================================
\section{The Electromagnetic Plane Wave}
%%=================================================================

\subsection{Vector potential, electric field, and magnetic field}

A monochromatic plane wave traveling in direction $\hat{\mathbf{q}}$
with polarization $\hat{\varepsilon}$ is described in the Coulomb
gauge ($\nabla\cdot\Avec=0$, $\phi=0$) by the vector potential
\begin{equation}
  \Avec(\mathbf{r},t)
  = A_0\,\hat{\varepsilon}\,
    e^{\iunit(\mathbf{q}\cdot\mathbf{r}-\omega t)}
    + \text{c.c.},
  \label{eq:A-planewave}
\end{equation}
where $\mathbf{q}$ is the photon wavevector, $|\mathbf{q}|=\omega/c$,
and $A_0$ is the field amplitude.

The Coulomb gauge condition $\nabla\cdot\Avec=0$ requires:
\begin{equation}
  \nabla\cdot\Avec
  = \iunit A_0\,(\hat{\varepsilon}\cdot\mathbf{q})\,
    e^{\iunit(\mathbf{q}\cdot\mathbf{r}-\omega t)} + \text{c.c.}
  = 0
  \quad\Rightarrow\quad
  \hat{\varepsilon}\cdot\mathbf{q} = 0.
  \label{eq:transversality}
\end{equation}
This is the \textbf{transversality condition}: the polarization
must be perpendicular to the propagation direction.
Electromagnetic waves in vacuum are transverse --- they cannot
have a field component along their direction of travel.

\medskip

The electric field is derived from $\Avec$ via
$\Evec = -\partial\Avec/\partial t$:
\begin{align}
  \Evec(\mathbf{r},t)
  &= -\frac{\partial\Avec}{\partial t}
   = \iunit\omega A_0\,\hat{\varepsilon}\,
     e^{\iunit(\mathbf{q}\cdot\mathbf{r}-\omega t)}
     + \text{c.c.}
  \label{eq:E-planewave}
\end{align}
So $\Evec$ oscillates in the \emph{same direction} as $\Avec$
(along $\hat{\varepsilon}$) and differs only by a factor of
$\iunit\omega$ (a time derivative introduces a factor of $-\iunit\omega$
acting on $e^{-\iunit\omega t}$, giving $+\iunit\omega$).

The magnetic field is $\Bvec = \nabla\times\Avec$.
Using $\nabla e^{\iunit\mathbf{q}\cdot\mathbf{r}}
= \iunit\mathbf{q}\,e^{\iunit\mathbf{q}\cdot\mathbf{r}}$
and the identity $\nabla\times(f\mathbf{v})
= (\nabla f)\times\mathbf{v}$ for scalar $f$, constant $\mathbf{v}$:
\begin{align}
  \Bvec(\mathbf{r},t)
  &= \nabla\times\Avec
   = \iunit\mathbf{q}\times
     \!\left(A_0\hat{\varepsilon}\right)
     e^{\iunit(\mathbf{q}\cdot\mathbf{r}-\omega t)}
     + \text{c.c.}
  \label{eq:B-planewave}
\end{align}
The direction of $\Bvec$ is along $\mathbf{q}\times\hat{\varepsilon}$,
perpendicular to both the propagation direction and the polarization.
The magnitudes satisfy $|\Bvec|/|\Evec| = |\mathbf{q}|/\omega = 1/c$
in Gaussian units (or $=1$ if we set $c=1$ in natural units),
confirming the standard equal-magnitude result for a plane wave.

\begin{keyresult}
\textbf{Summary:} For a plane wave propagating along $\hat{\mathbf{q}}$
with polarization $\hat{\varepsilon}$:
$\hat{\varepsilon}\perp\hat{\mathbf{q}}$ (transversality),
$\Evec\parallel\hat{\varepsilon}$,
$\Bvec\parallel\mathbf{q}\times\hat{\varepsilon}$,
and $\Avec$, $\Evec$, $\Bvec$ all carry the same spatial phase
$e^{\iunit\mathbf{q}\cdot\mathbf{r}}$.
\end{keyresult}

%%=================================================================
\section{Minimal Coupling: the Interaction Hamiltonian}
%%=================================================================

\subsection{Expanding the kinetic energy}

The single-electron Kohn-Sham Hamiltonian in the presence of the
EM field is obtained by minimal substitution
$\phat\to\phat+\Avec(\mathbf{r},t)$:
\begin{equation}
  \hat{H}(t)
  = \frac{1}{2}
    \!\left(\phat+\Avec(\mathbf{r},t)\right)^{\!2}
    + V_\mathrm{KS}(\mathbf{r},t).
  \label{eq:H-full}
\end{equation}
Expanding the square and using $\nabla\cdot\Avec=0$ (Coulomb gauge),
so that $\Avec\cdot\phat = \phat\cdot\Avec$:
\begin{equation}
  \hat{H}(t)
  = \underbrace{-\frac{\nabla^2}{2} + V_\mathrm{KS}}_{\hat{H}_0}
  + \underbrace{\Avec\cdot\phat}_{\hat{H}^{(1)}: \text{ linear in }\Avec}
  + \underbrace{\frac{1}{2}\Avec^2}_{\hat{H}^{(2)}: \text{ quadratic in }\Avec}.
  \label{eq:H-expand}
\end{equation}
The term $\hat{H}^{(2)} = \frac{1}{2}\Avec^2$ is the
\emph{diamagnetic} term; it is second order in the field amplitude
$A_0$ and is neglected in the linear response regime.
The linear interaction Hamiltonian is:
\begin{equation}
  \hat{H}_\mathrm{int}(\mathbf{r},t)
  = \Avec(\mathbf{r},t)\cdot\phat
  = A_0\,\hat{\varepsilon}\cdot\phat\;
    e^{\iunit(\mathbf{q}\cdot\mathbf{r}-\omega t)}
    + \text{c.c.}
  \label{eq:Hint-exact}
\end{equation}

%%=================================================================
\section{The Multipole Expansion}
%%=================================================================

\subsection{Expanding $e^{\iunit\mathbf{q}\cdot\mathbf{r}}$}

The spatial dependence of the interaction enters entirely through
$e^{\iunit\mathbf{q}\cdot\mathbf{r}}$ in eq.~(\ref{eq:Hint-exact}).
Taylor expanding:
\begin{equation}
  e^{\iunit\mathbf{q}\cdot\mathbf{r}}
  = 1
  + \iunit(\mathbf{q}\cdot\rhat)
  - \frac{1}{2}(\mathbf{q}\cdot\rhat)^2
  + \cdots
  \label{eq:Taylor}
\end{equation}
The matrix element of $\hat{H}_\mathrm{int}$ between orbitals
$|\psi_a\rangle$ and $|\psi_b\rangle$ becomes:
\begin{align}
  \langle\psi_b|\hat{H}_\mathrm{int}|\psi_a\rangle
  &= A_0 e^{-\iunit\omega t}
  \Big[
    \underbrace{
      \langle\psi_b|\hat{\varepsilon}\cdot\phat|\psi_a\rangle
    }_{\text{order 0: E1 (velocity)}}
    + \iunit
    \underbrace{
      \langle\psi_b|(\mathbf{q}\cdot\rhat)
      (\hat{\varepsilon}\cdot\phat)|\psi_a\rangle
    }_{\text{order 1: M1+E2}}
    + \cdots
  \Big]
  + \text{c.c.}
  \label{eq:multipole-me}
\end{align}

The successive terms are named after their physical content:

\begin{center}
\begin{tabular}{llll}
\hline
Term & Name & Operator & Size relative to E1 \\
\hline
Order 0 & E1: electric dipole (velocity form)
         & $\hat{\varepsilon}\cdot\phat$
         & 1 \\
Order 1 & M1: magnetic dipole\,+\,E2: electric quadrupole
         & $(\mathbf{q}\cdot\rhat)(\hat{\varepsilon}\cdot\phat)$
         & $|\mathbf{q}|\,a_0 \sim 10^{-3}$ \\
Order 2 & E3, M2, \ldots
         & $((\mathbf{q}\cdot\rhat)^2(\hat{\varepsilon}\cdot\phat)$
         & $(|\mathbf{q}|\,a_0)^2 \sim 10^{-6}$ \\
\hline
\end{tabular}
\end{center}

\subsection{Quantitative estimate for optical light}

For light at $\lambda = 400$~nm (near UV) in atomic units:
\begin{equation}
  |\mathbf{q}|
  = \frac{2\pi}{\lambda}
  = \frac{2\pi}{400\,\mathrm{nm}}
  = \frac{2\pi}{7559\,a_0}
  \approx 8.3\times10^{-4}\;a_0^{-1}.
\end{equation}
For a unit cell of linear size $a \sim 10$~\AA\ $= 18.9\,a_0$:
\begin{equation}
  |\mathbf{q}|\,a \approx 8.3\times10^{-4}\times18.9
  \approx 0.016 \ll 1.
\end{equation}
The M1+E2 corrections are therefore $\sim$1.6\% of E1 for a
10~\AA\ cell, and much smaller for typical bond lengths
($a_0$~scale).
The E1 approximation is excellent.

\subsection{The long-wavelength approximation}

Setting $e^{\iunit\mathbf{q}\cdot\mathbf{r}}\approx 1$ (keeping
only the order-0 term):
\begin{align}
  \Avec(\mathbf{r},t)
  &\xrightarrow{|\mathbf{q}|\,a\ll 1}
  \Avec(t)
  = A_0\,\hat{\varepsilon}\,e^{-\iunit\omega t}+\text{c.c.}
  \label{eq:A-uniform}
  \\[4pt]
  \hat{H}_\mathrm{int}(t)
  &= \Avec(t)\cdot\phat
  \label{eq:Hint-LWA}
\end{align}
The vector potential is now \emph{spatially uniform}: it depends
only on time, not on position.
Under this approximation, $\Bvec = \nabla\times\Avec(t) = \mathbf{0}$
--- the magnetic field associated with the uniform field is zero.
The perturbation is purely electric.

\begin{keyresult}
\textbf{Physical meaning of the long-wavelength approximation:}
The wavelength of the EM field is so large compared to the
system that every electron ``sees'' the same field at every
point in the unit cell.
The spatial phase $e^{\iunit\mathbf{q}\cdot\mathbf{r}}$ varies
negligibly across the cell and is replaced by~1.
This is equivalent to saying the photon wavevector is zero
($\mathbf{q}\to\mathbf{0}$).
\end{keyresult}

%%=================================================================
\section{Crystal Momentum Conservation and Vertical Transitions}
%%=================================================================

\subsection{General case: finite q transfers $\Delta\mathbf{k}=\mathbf{q}$}

For a periodic system with Bloch states
$\psi_{n\mathbf{k}}(\mathbf{r})
= e^{\iunit\mathbf{k}\cdot\mathbf{r}}u_{n\mathbf{k}}(\mathbf{r})$,
the matrix element of the full interaction~(\ref{eq:Hint-exact})
between states $|\psi_{n\mathbf{k}}\rangle$ and
$|\psi_{m\mathbf{k}'}\rangle$ involves the integral:
\begin{equation}
  \langle\psi_{m\mathbf{k}'}|
  e^{\iunit\mathbf{q}\cdot\mathbf{r}}\hat{\varepsilon}\cdot\phat
  |\psi_{n\mathbf{k}}\rangle
  \propto
  \int_\Omega\!
  e^{-\iunit\mathbf{k}'\cdot\mathbf{r}}\,u_{m\mathbf{k}'}^*\,
  e^{\iunit\mathbf{q}\cdot\mathbf{r}}\,
  e^{\iunit\mathbf{k}\cdot\mathbf{r}}\,
  u_{n\mathbf{k}}\,d^3r
  \propto \delta_{\mathbf{k}',\mathbf{k}+\mathbf{q}}.
  \label{eq:k-conservation}
\end{equation}
The delta function enforces \textbf{crystal momentum conservation}:
the photon transfers momentum $\mathbf{q}$ to the electron,
shifting its crystal momentum from $\mathbf{k}$ to
$\mathbf{k}+\mathbf{q}$.
The allowed transitions are:
\begin{equation}
  \Delta\mathbf{k} = \mathbf{q}_\mathrm{photon}.
  \label{eq:Deltak}
\end{equation}

\subsection{Long-wavelength limit: vertical transitions only}

Under the long-wavelength approximation $\mathbf{q}\to\mathbf{0}$:
\begin{equation}
  \langle\psi_{m\mathbf{k}'}|
  \Avec(t)\cdot\phat
  |\psi_{n\mathbf{k}}\rangle
  \propto \delta_{\mathbf{k}',\mathbf{k}}.
  \label{eq:vertical}
\end{equation}
All transitions are \textbf{vertical} ($\Delta\mathbf{k}=0$) in
the Brillouin zone --- the optical response probes only
interband transitions at the same $\mathbf{k}$-point.
This is the standard result for optical spectroscopy and is the
regime of paper~2.

\subsection{Finite-q response: beyond vertical transitions}

If the long-wavelength approximation is relaxed and $\mathbf{q}$
is treated as a free parameter (decoupled from $\omega/c$),
one can directly compute the \textbf{density response function}
$\chi(\mathbf{q},\omega)$ by applying a density-wave perturbation:
\begin{equation}
  \delta V(\mathbf{r},t)
  = V_0\,e^{\iunit\mathbf{q}\cdot\mathbf{r}}\,\delta(t).
  \label{eq:density-wave-kick}
\end{equation}
This is a scalar potential perturbation (length-gauge language)
with the spatial profile $e^{\iunit\mathbf{q}\cdot\mathbf{r}}$.
The time-domain response is:
\begin{equation}
  \delta\rho(\mathbf{q},t)
  = \int_\Omega\!
    \delta\rho(\mathbf{r},t)\,
    e^{-\iunit\mathbf{q}\cdot\mathbf{r}}\,d^3r,
\end{equation}
and the Fourier transform in time gives:
\begin{equation}
  \chi(\mathbf{q},\omega)
  = \frac{\widetilde{\delta\rho}(\mathbf{q},\omega)}{V_0},
  \qquad
  S(\mathbf{q},\omega) \propto \mathrm{Im}\,\chi(\mathbf{q},\omega).
  \label{eq:chi-and-S}
\end{equation}
Repeating for a sequence of commensurate $\mathbf{q}$ values
maps out the dynamic structure factor $S(\mathbf{q},\omega)$,
giving access to plasmon dispersion, phonon branches (via the
electron-density response to lattice distortions), and
interband transitions at finite momentum transfer.
The constraint is that $\mathbf{q}$ must be commensurate with
the simulation cell:
\begin{equation}
  \mathbf{q}
  = \frac{n_1}{N_1}\mathbf{b}_1
  + \frac{n_2}{N_2}\mathbf{b}_2
  + \frac{n_3}{N_3}\mathbf{b}_3,
  \label{eq:q-commensurate}
\end{equation}
where $\mathbf{b}_i$ are reciprocal lattice vectors and $N_i$
are the supercell repeat dimensions.

\begin{keyresult}
\textbf{Key distinction:}
The long-wavelength approximation sets $\mathbf{q}=\mathbf{0}$
and recovers the optical response (vertical transitions).
The finite-$\mathbf{q}$ generalization~(\ref{eq:density-wave-kick})
is a length-gauge scalar potential perturbation that computes
$S(\mathbf{q},\omega)$, comparing directly with inelastic
X-ray scattering and EELS experiments.
The photon coupling to $\Avec$ (momentum matrix) and the density
response to $\delta V$ (density-wave kick) are related but
distinct formalisms; the connection is through the fluctuation-dissipation
theorem and the Ward identity of the linear response.
\end{keyresult}

%%=================================================================
\section{Connection to the Coulomb-Matrix Perturbation}
%%=================================================================

An alternative perturbation that gives spatially resolved (real-space)
response is the Coulomb potential of a point charge placed at
position $\mathbf{R}$ inside the simulation cell:
\begin{equation}
  \delta V_\mathrm{Coul}(\mathbf{r})
  = \frac{V_0}{|\mathbf{R}-\mathbf{r}|},
  \label{eq:Coulomb-pert}
\end{equation}
with matrix elements
$\langle\phi_i|1/|\mathbf{R}-\hat{\mathbf{r}}||\phi_j\rangle$.
This approach, implemented in ref.~[Jakowski et al.\ JCTC 2019],
mimics a focused electron beam at position $\mathbf{R}$.

The connection to the plane-wave perturbation~(\ref{eq:density-wave-kick})
is through the Fourier expansion of the Coulomb kernel:
\begin{equation}
  \frac{1}{|\mathbf{R}-\mathbf{r}|}
  = \frac{4\pi}{\Omega}
    \sum_{\mathbf{q}\neq\mathbf{0}}
    \frac{e^{\iunit\mathbf{q}\cdot(\mathbf{r}-\mathbf{R})}}
         {|\mathbf{q}|^2},
  \label{eq:Coulomb-Fourier}
\end{equation}
where the sum is over reciprocal lattice vectors $\mathbf{q}$ of
the simulation cell.
Comparing eqs.~(\ref{eq:density-wave-kick}) and (\ref{eq:Coulomb-Fourier}):

\begin{keyresult}
The Coulomb perturbation at position $\mathbf{R}$ is a
\textbf{superposition of all plane-wave perturbations
$e^{\iunit\mathbf{q}\cdot\mathbf{r}}$ weighted by
$e^{-\iunit\mathbf{q}\cdot\mathbf{R}}/|\mathbf{q}|^2$}.
\begin{itemize}
  \item The $1/|\mathbf{q}|^2$ weight emphasizes small-$\mathbf{q}$
    (long-wavelength) components, reflecting the long-range nature
    of the Coulomb interaction.
  \item The Coulomb approach gives the response as a function of
    probe position $\mathbf{R}$ (real space) --- ideal for simulating
    a focused electron beam in STEM-EELS.
  \item The plane-wave approach gives the response as a function of
    $\mathbf{q}$ (reciprocal space) --- ideal for plasmon dispersion
    and comparison with momentum-resolved EELS or inelastic X-ray
    scattering.
  \item They contain the same physical information and are
    related by a Fourier transform; neither is more fundamental.
\end{itemize}
\end{keyresult}

%%=================================================================
\section{The Length Gauge Hamiltonian}
%%=================================================================

\subsection{Length gauge with the full plane wave (finite q)}

In the length gauge (LG), the EM interaction is described entirely
through the scalar potential ($\Avec^\mathrm{LG}=\mathbf{0}$,
$\phi^\mathrm{LG}=-\mathbf{r}\cdot\mathbf{E}$).
For a spatially varying field~--- keeping the full
$e^{\iunit\mathbf{q}\cdot\mathbf{r}}$ factor --- the scalar
potential is:
\begin{equation}
  \phi^\mathrm{LG}(\mathbf{r},t)
  = -\frac{1}{|\mathbf{q}|}
    E_0\,e^{\iunit(\mathbf{q}\cdot\mathbf{r}-\omega t)}
    + \text{c.c.},
  \label{eq:phi-LG-finiteq}
\end{equation}
and the interaction Hamiltonian is:
\begin{equation}
  \hat{H}_\mathrm{int}^\mathrm{LG}(\mathbf{r},t)
  = -\phi^\mathrm{LG}(\mathbf{r},t)
  = \frac{E_0}{|\mathbf{q}|}
    e^{\iunit(\mathbf{q}\cdot\mathbf{r}-\omega t)}
    + \text{c.c.}
  \label{eq:HLG-finiteq}
\end{equation}
The matrix element between Bloch states is:
\begin{equation}
  \langle\psi_{m\mathbf{k}'}|
  \hat{H}_\mathrm{int}^\mathrm{LG}
  |\psi_{n\mathbf{k}}\rangle
  \propto
  \langle u_{m\mathbf{k}}|e^{\iunit\mathbf{q}\cdot\mathbf{r}}|u_{n\mathbf{k}}\rangle
  \cdot\delta_{\mathbf{k}',\mathbf{k}+\mathbf{q}}.
  \label{eq:HLG-me}
\end{equation}
For $\mathbf{q}\to\mathbf{0}$, $e^{\iunit\mathbf{q}\cdot\mathbf{r}}\to 1$
and this approaches the dipole matrix element
$\langle u_{m\mathbf{k}}|u_{n\mathbf{k}}\rangle = \delta_{mn}$
for orthogonal states, which vanishes --- the coupling comes from
the $\mathbf{q}\cdot\mathbf{r}$ term and gives
$\langle u_{m\mathbf{k}}|\hat{r}_\alpha|u_{n\mathbf{k}}\rangle$:
the standard length-gauge dipole matrix element of paper~1.

\subsection{Why the length gauge fails for periodic systems at q=0}

As noted above, for $\mathbf{q}\to\mathbf{0}$ the length gauge
requires matrix elements of $\hat{\mathbf{r}}$ between periodic
Bloch states.
These matrix elements are ill-defined for extended states:
$\langle\psi_{n\mathbf{k}}|\hat{r}_\alpha|\psi_{m\mathbf{k}}\rangle$
is not translationally invariant and diverges for macroscopic
systems.
This is the fundamental reason the velocity gauge is required for
periodic RT-TDDFT.

For \emph{finite-$\mathbf{q}$} perturbations in the length gauge,
the ill-defined $\hat{\mathbf{r}}$ does NOT appear: the matrix
element involves $e^{\iunit\mathbf{q}\cdot\mathbf{r}}$, which is
periodic and well-defined for any $\mathbf{q}\neq\mathbf{0}$.
The length-gauge difficulty is specific to $\mathbf{q}=0$ only.

%%=================================================================
\section{The Velocity Gauge Hamiltonian}
%%=================================================================

\subsection{Gauge function and transformation rules}

The length and velocity gauges are connected by the gauge function:
\begin{equation}
  \chi(\mathbf{r},t)
  = -\hat{\mathbf{r}}\cdot\Avec^\mathrm{VG}(t),
  \label{eq:chi}
\end{equation}
where $\Avec^\mathrm{VG}(t) = -\int_{-\infty}^t\mathbf{E}(t')dt'$
is the velocity-gauge vector potential.
Under this gauge transformation, the potentials, wavefunctions,
and Hamiltonians transform as:
\begin{align}
  \Avec &\to \Avec' = \Avec + \nabla\chi,
  \label{eq:gt-A}
  \\
  \phi &\to \phi' = \phi - \frac{\partial\chi}{\partial t},
  \label{eq:gt-phi}
  \\
  \Psi &\to \Psi' = \hat{U}\Psi,
  \quad
  \hat{U} = e^{-\iunit\chi} = e^{\iunit\mathbf{r}\cdot\Avec(t)},
  \label{eq:gt-psi}
  \\
  \hat{H} &\to \hat{H}'
  = \hat{U}\hat{H}\hat{U}^\dagger
    - \iunit\hat{U}\frac{\partial\hat{U}^\dagger}{\partial t}.
  \label{eq:gt-H}
\end{align}
The unitary operator $\hat{U}$ preserves the charge density:
$|\Psi'|^2 = |\Psi|^2$.

\subsection{Deriving H$^\mathrm{VG}$ step by step}

We start from the length-gauge Hamiltonian:
\begin{equation}
  \HLG(t)
  = -\frac{\nabla^2}{2} + V_\mathrm{KS}(\mathbf{r},t)
  + \Vnl
  + \hat{\mathbf{r}}\cdot\mathbf{E}(t),
  \label{eq:HLG}
\end{equation}
and apply the transformation~(\ref{eq:gt-H}).

\paragraph{Step 1: Transform the kinetic energy.}
We need $\hat{U}\phat_\alpha\hat{U}^\dagger$.
Acting on a test function $g(\mathbf{r})$:
\begin{align}
  \hat{U}\phat_\alpha\hat{U}^\dagger g
  &= e^{\iunit\mathbf{r}\cdot\Avec}
     (-\iunit\nabla_\alpha)
     \!\left(e^{-\iunit\mathbf{r}\cdot\Avec}g\right)
  \nonumber\\
  &= e^{\iunit\mathbf{r}\cdot\Avec}
     (-\iunit)
     \!\left[(-\iunit A_\alpha)
     e^{-\iunit\mathbf{r}\cdot\Avec}g
     + e^{-\iunit\mathbf{r}\cdot\Avec}(\nabla_\alpha g)
     \right]
  \nonumber\\
  &= (\hat{p}_\alpha - A_\alpha)\,g.
  \label{eq:Uphat}
\end{align}
Therefore $\hat{U}\phat\hat{U}^\dagger = \phat - \Avec$, and:
\begin{equation}
  \hat{U}\!\left(-\frac{\nabla^2}{2}\right)\hat{U}^\dagger
  = \frac{1}{2}(\phat - \Avec)^2.
  \label{eq:UT}
\end{equation}
For an electron (charge $-1$), the LG kinetic term already has the
sign convention $+\frac{1}{2}\phat^2$; after transformation this
becomes $\frac{1}{2}(\phat-\Avec)^2$.
Accounting for the electron charge sign throughout
(which appears as $+\Avec$ in the VG Hamiltonian --- see footnote
in Section 2.3 of the manuscript):
\begin{equation}
  \text{kinetic term:}\quad
  \frac{1}{2}\hat{p}^2
  \xrightarrow{\hat{U}}
  \frac{1}{2}(\phat+\Avec)^2.
  \label{eq:kinetic-VG}
\end{equation}

\paragraph{Step 2: Transform the time-derivative term.}
\begin{align}
  -\iunit\hat{U}\frac{\partial\hat{U}^\dagger}{\partial t}
  &= -\iunit\,e^{\iunit\mathbf{r}\cdot\Avec}
     \frac{\partial}{\partial t}
     e^{-\iunit\mathbf{r}\cdot\Avec}
  \nonumber\\
  &= -\iunit\,e^{\iunit\mathbf{r}\cdot\Avec}
     \cdot(-\iunit\,\mathbf{r}\cdot\dot{\Avec})
     e^{-\iunit\mathbf{r}\cdot\Avec}
  \nonumber\\
  &= -\,\mathbf{r}\cdot\dot{\Avec}(t)
  = +\,\mathbf{r}\cdot\mathbf{E}(t),
  \label{eq:dt-term}
\end{align}
where we used $\mathbf{E} = -\partial\Avec/\partial t$.
This term is $+\mathbf{r}\cdot\mathbf{E}$, which exactly cancels
the length-gauge coupling $+\mathbf{r}\cdot\mathbf{E}$ in
eq.~(\ref{eq:HLG}):
\begin{equation}
  \underbrace{-\mathbf{r}\cdot\mathbf{E}}_{\text{from }\hat{U}\HLG\hat{U}^\dagger}
  + \underbrace{+\mathbf{r}\cdot\mathbf{E}}_{\text{from }-\iunit\hat{U}\partial_t\hat{U}^\dagger}
  = 0.
  \label{eq:LG-coupling-cancels}
\end{equation}
This cancellation is the central result of the Göppert-Mayer
transformation: the length-gauge interaction $\mathbf{r}\cdot\mathbf{E}$
disappears entirely.

\paragraph{Step 3: Transform the local potential.}
$V_\mathrm{KS}(\mathbf{r})$ is multiplicative, so it commutes with
$\hat{U}(\mathbf{r},t)$:
\begin{equation}
  \hat{U}V_\mathrm{KS}\hat{U}^\dagger = V_\mathrm{KS}.
  \label{eq:VKS-invariant}
\end{equation}

\paragraph{Step 4: Transform the nonlocal pseudopotential.}
$\Vnl$ is nonlocal (not diagonal in position space), so
$[\hat{U},\Vnl]\neq0$.
In the linear-response regime ($A_0\ll 1$), expand:
\begin{align}
  \hat{U}\Vnl\hat{U}^\dagger
  &= e^{\iunit\mathbf{r}\cdot\Avec}\Vnl e^{-\iunit\mathbf{r}\cdot\Avec}
  \nonumber\\
  &\approx \Vnl
  + \iunit\Avec(t)\cdot[\hat{\mathbf{r}},\Vnl]
  + \mathcal{O}(A_0^2).
  \label{eq:VNL-transform}
\end{align}
The correction $\iunit\Avec\cdot[\hat{\mathbf{r}},\Vnl]$ is the
gauge-transformation contribution of the nonlocal pseudopotential.
As shown in the manuscript appendix, this is identical to the
Starace commutator that appears in the velocity operator:
$\hat{v}_\alpha = \hat{p}_\alpha + \frac{1}{\iunit}[\hat{r}_\alpha,\Vnl]$.

\paragraph{Collecting all terms:}
\begin{align}
  \HVG(t)
  &= \hat{U}\HLG\hat{U}^\dagger
   - \iunit\hat{U}\partial_t\hat{U}^\dagger
  \nonumber\\
  &= \underbrace{\frac{1}{2}(\phat+\Avec)^2}_{\text{Step 1}}
   + \underbrace{V_\mathrm{KS}}_{\text{Step 3}}
   + \underbrace{\Vnl
     + \iunit\Avec\cdot[\hat{\mathbf{r}},\Vnl]}_{\text{Step 4}}
   + \underbrace{\cancel{\mathbf{r}\cdot\mathbf{E}}
     - \cancel{\mathbf{r}\cdot\mathbf{E}}}_{\text{Step 2: cancels}}
   \nonumber\\[6pt]
  &\boxed{
   = \frac{1}{2}(\phat+\Avec)^2
   + V_\mathrm{KS}(\mathbf{r},t)
   + \Vnl
   + \iunit\Avec(t)\cdot[\hat{\mathbf{r}},\Vnl].
   }
  \label{eq:HVG-final}
\end{align}

\begin{keyresult}
\textbf{Velocity-gauge Kohn-Sham Hamiltonian:}
\begin{equation*}
  \HVG(t)
  = \frac{1}{2}\!\left(\hat{\mathbf{p}} + \Avec(t)\right)^{\!2}
  + V_\mathrm{KS}(\mathbf{r},t)
  + \Vnl
  + \iunit\Avec(t)\cdot\!\left[\hat{\mathbf{r}},\Vnl\right]
\end{equation*}
\textbf{Key features:}
\begin{itemize}
  \item No $\hat{\mathbf{r}}$ operator anywhere --- valid for
    periodic systems.
  \item EM coupling appears through $\Avec(t)\cdot\phat$ (linear)
    and $\Avec^2/2$ (diamagnetic, dropped in linear response).
  \item Nonlocal PP acquires a correction
    $\iunit\Avec\cdot[\hat{\mathbf{r}},\Vnl]$, equivalent to the
    Starace term in the current operator.
  \item For spatially uniform $\Avec(t)$ (long-wavelength limit),
    all $\mathbf{k}$-points experience the same perturbation.
\end{itemize}
\end{keyresult}

%%=================================================================
\section{The Delta-Kick Perturbation in Velocity Gauge}
%%=================================================================

In linear response, the vector potential is a delta-function impulse:
\begin{equation}
  \Avec(t) = A_0\hat{\varepsilon}\,\delta(t).
  \label{eq:VG-kick}
\end{equation}
At $t=0^-$ the system is in the ground state $\psi_{n\mathbf{k}}^{(0)}$.
Integrating the TDKS equation through the delta function,
the orbitals acquire a phase at $t=0^+$:
\begin{equation}
  \psi_{n\mathbf{k}}(\mathbf{r},0^+)
  = e^{\iunit A_0\hat{\varepsilon}\cdot\phat}\,
    \psi_{n\mathbf{k}}^{(0)}(\mathbf{r}).
  \label{eq:kick-phase}
\end{equation}
For small $|A_0|\ll 1$ (linear response):
\begin{equation}
  \psi_{n\mathbf{k}}(\mathbf{r},0^+)
  \approx \left(1 + \iunit A_0\hat{\varepsilon}\cdot\phat\right)
  \psi_{n\mathbf{k}}^{(0)}(\mathbf{r}).
  \label{eq:kick-linear}
\end{equation}
For $t>0$: $\Avec(t)=0$ and the system evolves under
$\hat{H}_0 + V_\mathrm{KS}[\rho(t)]$ alone.
The current $J_\alpha(t)$ oscillates in response to the kick and
encodes the full optical response via Fourier transformation.

%%=================================================================
\section{Summary: Hierarchy of Approximations}
%%=================================================================

\begin{center}
\begin{tabular}{p{3.2cm}p{4.8cm}p{5cm}}
\hline
\textbf{Approximation} & \textbf{What is kept} & \textbf{Physical content} \\
\hline
Full plane wave, LG
& $e^{\iunit\mathbf{q}\cdot\mathbf{r}}\delta(t)$,
  scalar potential
& $\chi(\mathbf{q},\omega)$, $S(\mathbf{q},\omega)$,
  EELS, IXS, plasmon dispersion \\[4pt]
Long-wavelength LG
& $\mathbf{r}\cdot\mathbf{E}(t)$ ($\mathbf{q}\to0$)
& Optical response, vertical transitions;
  ill-defined for PBC \\[4pt]
Long-wavelength VG
& $\Avec(t)\cdot\phat$ ($\mathbf{q}\to0$)
& Same optical response; well-defined for PBC.
  \textbf{This paper.} \\[4pt]
Delta-kick VG
& $A_0\hat{\varepsilon}\cdot\phat\,\delta(t)$
& Full spectrum in one run;
  $\sigma_{\alpha\beta}(\omega)$ from FFT \\
\hline
\end{tabular}
\end{center}

\vspace{12pt}

The Coulomb-matrix perturbation
$\langle\phi_i|1/|\mathbf{R}-\hat{\mathbf{r}}||\phi_j\rangle\,\delta(t)$
is a real-space version of the ``full plane wave, LG'' row,
with $1/|\mathbf{q}|^2$ weighting across all $\mathbf{q}$.
It gives spatially resolved response (STEM-EELS analogue) while
the pure $e^{\iunit\mathbf{q}\cdot\mathbf{r}}$ kick gives
momentum-resolved response --- complementary descriptions of the
same underlying density response function.

\end{document}
